%%%%%%%%%%  scan idx 89  |  folio 76  %%%%%%%%%%
De la
formule (7) résulte aussitôt \(y\lambda_{-1}(\check F)\in j^!(K(X'))\) pourvu que l'on ait~:
\begin{equation}\tag{11}
  y\,\gamma^{p-1}(E)\lambda^p(\check E)\in i^!(K(X)).
\end{equation}
Supposons cette condition (11) satisfaite. D'après le corollaire au lemme 2, \(\RR\) est vrai pour \((Y',X',y\lambda_{-1}(\check F))\). D'autre part, pour prouver (2), il suffit, en vertu de la formule (8), de prouver que
\begin{equation}\tag{12}
  f^*\ch(i_!(y))
  \equiv
  f^*\bigl(i_*(\ch(y)T(E)^{-1})\bigr)
  \quad\text{mod.}\Ker f_*.
\end{equation}
Or, cette dernière formule se vérifie aussitôt, utilisant (6), (9), \(\RR\) pour \((Y',X',y\lambda_{-1}(F))\) et la formule :
\begin{equation}\tag{13}
  \ch(\lambda_{-1}(\check F))=C^{p-1}(F)T(F)^{-1}\quad\text{(cf.\ formule (5)).}
\end{equation}

Ainsi, \(\RR\) est prouvé chaque fois que \(y\,\gamma^{p-1}(E)=0\), en particulier (puisque \(\gamma^{p-1}(E)\) est de filtration \(\geq p-1\) {*}dans \(K(Y)\) filtré par la codimension des supports{*}) chaque fois que \(p-1>\dim.Y\), i.e.\ chaque fois que \(\dim.Y<\frac{\dim.X-1}{2}\). Or, on se ramène à ce cas, en conjuguant le lemme 1 et le lemme 3 (prenant pour \(P\) un espace projectif de dimension \(r\) assez grande, par exemple \(r>\dim.X+1\)). C.Q.F.D.

N.B. On s'est servi du fait que \(\gamma^{p-1}(E)\) était de filtration \(\geq p-1\). On peut prouver ce fait directement, sans passer par les grassmaniennes comme dans \(\mathrm{R.R.R.}\) 2.16, par l'étude du fibré \(Y'\) en espaces projectifs associé à \(E\). Cette méthode donne plus généralement le th.\ 2.14 de \(\mathrm{R.R.R.}\)

\subsection*{Extrait d'une lettre de Grothendieck}
\dots{} «~moralement~», la nouvelle méthode repose sur la détermination de \(K(X)\) et de \(A(X)\) lorsque \(X\) est, soit un fibré projectif associé à un fibré vectoriel, soit une variété obtenu par éclatement d'une sous-variété non
